\section{Conclusão}

Neste trabalho foi estudado o Problema do Caminho Hamiltoniano de forma teórica e experimental. Inicialmente, foi apresentada a definição formal do problema e provada formalmente o seu pertencimento à classe NP e sua NP-completude por meio de reduções, começando a partir do 3-SAT para o problema do Caminho Hamiltoniano Dirigido, que por sua vez foi reduzido para o caso não direcionado, que é o foco deste trabalho. 

Também foram discutidos, de forma mais introdutória, métodos clássicos como backtracking e branch-and-bound, bem como heurísticas (como métodos gulosos, método de Pósa e metaheurísticas) e aproximações propostas na literatura, que embora mais rápidas não garantem exatidão. Essas abordagens foram apresentadas principalmente para contextualização. Em seguida, foram analisadas duas abordagens exatas para sua resolução: um algoritmo baseado em programação dinâmica (BMK) e uma abordagem baseada em redução para SAT resolvida por um solucionador CDCL (RedSAT).

Os resultados dos experimentos mostraram que o  BMK apresenta desempenho superior em termos de tempo de execução para todas as instâncias consideradas. Por outro lado, a abordagem baseada em SAT apresentou crescimento significativamente mais rápido do tempo de execução.

Em termos de uso de memória, ambas as abordagens apresentaram consumo moderado para os tamanhos considerados no experimento. Entretanto, a análise teórica indica que o consumo de memória do BMK explode a partir de valores de $n$ um pouco maiores do que os deste trabalho. Já o uso de memória no RedSAT tende a crescer de forma mais suave em relação ao BMK, ainda que o consumo seja relativamente bem alto para valores de $n$ menores.

Os resultados indicam que, para instâncias pequenas e médias do Problema do Caminho Hamiltoniano, um algoritmo exato especializado é preferível a uma formulação genérica via SAT, tanto por razões de eficiência quanto de simplicidade de implementação. No entanto, para valores de $n$ suficientemente grandes ambas abordagens são computacionalmente inviáveis devido às suas complexidades exponenciais.

Como trabalhos futuros, podem ser analisados o uso de heurísticas e algoritmos aproximados para instâncias maiores, onde também poem ser analisada a precisão dos mesmos (afinal, verificar a corretude da resposta é polinomial). Também podem ser investigados métodos híbridos, que combinem métodos exatos (como programação dinâmica) com heurísticas e/ou aproximações.